\begin{figure}[H]
    \centering
    \begin{tikzpicture}[x=1.0cm, y=1.0cm]

        % (a) the threshold test
        \node[panel, font=\scriptsize\bfseries] at (-0.55,1.55) {(a) the root threshold, $k=4$};

        \node[blkreg, minimum width=0.66cm] (r)  at (1.55,0.85) {$6$};
        \node[blkreg, minimum width=0.66cm] (c1) at (0.85,-0.05) {$8$};
        \node[blkreg, minimum width=0.66cm] (c2) at (2.25,-0.05) {$7$};
        \node[blkreg, minimum width=0.66cm] (c3) at (0.15,-0.95) {$9$};
        \draw[flowmute, -] (r) -- (c1);
        \draw[flowmute, -] (r) -- (c2);
        \draw[flowmute, -] (c1) -- (c3);
        \node[notemute, font=\tiny, anchor=east] at (1.18,0.85) {root $=$ threshold};
        \node[notemute, font=\tiny, anchor=north, align=center] at (1.55,-1.42)
            {min-heap with the\\ k best (so far)};

        \node[lanedrop] (x1) at (4.30,0.85) {$3$};
        \node[note, font=\tiny, anchor=west] at (4.68,0.85) {$3 \le 6$};
        \draw[flowmute, dashed, dash pattern=on 2pt off 1.6pt] (x1) -- (r);
        \node[note, font=\scriptsize, anchor=west, text=figregD, align=left] at (6.30,0.85)
            {reject:\\ \emph{one} comparison};

        \node[lanereg] (x2) at (4.30,-0.95) {$11$};
        \node[note, font=\tiny, anchor=west] at (4.68,-0.95) {$11 > 6$};
        \draw[flow] (x2) -- (r);
        \node[note, font=\scriptsize, anchor=west, text=figdrmD, align=left] at (6.30,-0.95)
            {replace root\\ and sift down: $O(\log k)$};

        % (b) map-reduce
        \begin{scope}[yshift=-3.55cm]
            \node[panel, font=\scriptsize\bfseries] at (-0.55,1.10) {(b) the map-reduce scheme};

            \foreach \i in {0,1,2,3} {
                \pgfmathsetmacro{\xc}{1.05+\i*1.85}
                \node[blk, minimum width=1.55cm, minimum height=0.40cm] (ch\i) at (\xc,0.52)
                    {\tiny chunk \i};
                \node[blkreg, minimum width=0.85cm, minimum height=0.38cm] (hp\i) at (\xc,-0.42)
                    {\tiny heap k};
                \draw[flow] (ch\i.south) -- (hp\i.north);
            }
            \node[rowlab, font=\tiny] at (ch0.west) {map};

            \node[blkreg, minimum width=1.55cm, minimum height=0.40cm] (mg) at (3.83,-1.45)
                {\tiny final k};
            \foreach \i in {0,1,2,3} { \draw[flowmute] (hp\i.south) -- (mg.north); }
            \node[rowlab, font=\tiny] at (0.10,-1.45) {reduce};

            \node[note, font=\scriptsize, anchor=west, align=left] at (7.70,-0.42)
                {no communication\\ in the map phase};
            \node[note, font=\scriptsize, anchor=west, align=left, text=figdrmD] at (7.70,-1.45)
                {each heap filters only\\ after seeing k elements};
        \end{scope}

    \end{tikzpicture}
    \caption[Heap-based selection and the map-reduce scheme]{Heap-based selection and its parallelisation. In (a) the root of the
    min-heap acts as a threshold: an element that does not exceed it is
    rejected with one comparison, while one that exceeds it replaces the root
    and triggers a sift-down. In (b) the input is partitioned into chunks, one
    heap per worker, and the partial results are merged into one.}
    \label{fig:bg-heap}
\end{figure}
